\section{Introduction}

Cross-market replication is most informative when it tests not only whether a return
pattern survives, but whether the state variable retains the economic meaning assigned
to it in the source market. This distinction matters for volatility-based strategy
selection. In equities, high market volatility often co-occurs with market distress,
funding constraints, liquidity demand, and subsequent rebound states in which momentum
can crash and recent losers can reverse. A scalar volatility measure can therefore act
as a compact proxy for a richer joint state.

\citet{butt2023} use this idea to propose a constant-leverage switching rule. The rule
holds cross-sectional momentum when current market volatility lies in the first three
quartiles of its trailing five-year distribution and holds short-term reversal in the
highest quartile. Its appeal is practical as well as economic: it avoids the ex-post
leverage normalization used by some volatility-scaled momentum strategies and can be
specified before next-month returns are observed.

We ask whether that same state variable is sufficient in Chinese futures. Futures are a
demanding transportability test. Their cross-section spans financial contracts,
industrial inputs, agricultural commodities, energy, and other products whose price
changes originate in different cash-market constraints. They also lack a unique analogue
of the equity market-capitalization portfolio. A high aggregate volatility observation
may reflect a broad deleveraging episode, a sector-specific supply shock, a policy change,
or several unrelated shocks occurring simultaneously. These states need not imply the
same next-month ranking of momentum and reversal.

The empirical design is deliberately restrictive. We translate the published rule before
viewing the replication returns; retain the fixed 75th-percentile boundary; use all
eligible products rather than an outcome-selected subset; and report two defensible
futures-market return proxies separately. We do not search for a better percentile after
observing the results. This converts proxy sensitivity and event concentration into
falsification evidence rather than tuning opportunities.

The results do not support a robust replication. The primary switch earns a gross
arithmetic annual return of 13.69\%, compared with 13.25\% for momentum, but the
increment has a HAC $t$-statistic of only 0.06. Reversal beats momentum in 4 of 14
primary-proxy high-volatility months. The mean advantage is positive only because of one
large event; the median and the mean excluding that event are negative. More importantly,
an equal-product volatility proxy reverses the Q4 sleeve ranking and produces a router
that underperforms momentum. The two volatility series have a correlation of 0.83, yet
their quartile states agree in only 42.4\% of strategy months.

The negative result is not equivalent to saying that volatility is noise. In the primary
proxy, Q4 signal months have the highest average pairwise product correlation,
cross-sectional dispersion, and directional coherence. Volatility successfully measures
the intensity and commonality of shocks. What it fails to identify is whether the shock
is trend-generating or reversal-generating.

The paper makes four contributions. First, it provides a transparent cross-asset
replication of a fixed equity routing rule in a heterogeneous futures market. Second, it
shows why high correlation between alternative state measures is not enough when the
decision depends on crossing a quantile boundary. Third, it separates an event-level
crash-hedge echo from a stable conditional relation. Fourth, it develops an economic
interpretation of state insufficiency: volatility is an intensity scalar, whereas
strategy selection requires information about direction, source, breadth, and shock
stage.
